\documentclass{ximera}


% MATH -----------------------------------------------------------
\newcommand{\nats}{\mathbb N}
\newcommand{\ints}{\mathbb Z}
\newcommand{\rats}{\mathbb Q}
\newcommand{\reals}{\mathbb R}
\newcommand{\complex}{\mathbb C}
\newcommand{\powerset}{\mathscr P}

%-------------------------------------------------------------

\pgfplotsset{compat=1.18}
\begin{document}
\begin{abstract}
 \end{abstract}

    \title{Class Discussion Activity 10}
    
    \maketitle

\section*{Exercises}



\begin{enumerate}[label=(\arabic*)]

\item Suppose the following incorrect work is a student's attempt to use Variation of Parameters to find a solution to $(x-1)y^{\,\prime\prime}-xy^{\,\prime}+y=2(x-1)^2\,e^{x}$ given that $y_1=x$ and $y_2=e^x$ form a fundamental set of solutions to the complementary equation.  What error is this student making?

\emph{The Wronskian of $y_1=x$ and $y_2=e^x$ is $W=y_1y_2^{\,\prime}-y_1^{\prime}y_2=(x-1)e^x$.  Then there is a particular solution to $(x-1)y^{\,\prime\prime}-xy^{\,\prime}+y=2(x-1)^2\,e^{x}$ of the form $y_p=uy_1+vy_2$ where $\displaystyle u^{\,\prime}=-\frac{2(x-1)^2e^{x}e^x}{(x-1)e^x}=-2(x-1)e^{x}$ and \\ $\displaystyle v^{\,\prime}=\frac{2(x-1)^2e^{x}x}{(x-1)e^x}=2x^2-2x$.  By integrating, we get $u=-2(x-2)e^x$ and $\displaystyle v=\frac{2}{3}x^3-x^2$.  So a solution is \\$y_p=ux+ve^x=(-2(x-2)e^x)x+\left(\frac{2}{3}x^3-x^2\right)e^x=\left(\frac{2}{3}x^3-3x^2+4x\right)e^x$. }


\begin{enumerate}[label=(\alph*)]

\item Variation of Parameters cannot be used to find a solution to this differential equation.

\item The Wronskian of $y_1=x$ and $y_2=e^x$ is not $W=(x-1)e^x$.

\item An incorrect ``forcing" function $f(x)$ is being used in the formulas for $u^{\,\prime}$ and $v^{\,\prime}$.  % ***

\item The function $-2(x-2)e^x$ is not an antiderivative of $-2(x-1)e^{x}$.

\item The function $\frac{2}{3}x^3-x^2$ is not an antiderivative of $2x^2-2x$.

\item $(-2(x-2)e^x)x+\left(\frac{2}{3}x^3-x^2\right)e^x \neq \left(\frac{2}{3}x^3-3x^2+4x\right)e^x$.

\end{enumerate}

% (c) with options through (f)

\item Correct the error above and determine the correct particular solution $y_p$ that should have been determined.  Determine $y_p(3)$.  Round to the nearest hundredth.

% 3e^3 \approx 60.26

\newpage


\item Suppose Variation of Parameters was used to find the solution to \\ $y^{\,\prime\prime}-y=f(t)$ and it was found that $u'=3t^2+3t$ and $v'=-(3t^2+3t)e^{2t}$ using complementary solutions $y_1=e^t$ and $y_2=e^{-t}$.  Find $f(1)$ to the nearest hundredth.

% Say a solution is y=t^3e^t.  y'=(t^3+3t^2)e^t and y''=(t^3+6t^2+6t)e^t

% So f(t)=(6t^2+6t)e^t

% f(1)=12e\approx 32.62

%  y_1=e^t, y_2=e^(-t), f=(6t^2+6t)e^t

% W=-2

% u'=-fy_2/W=(6t^2+6t)/2  --> u'=3t^2+3t ---> u=t^3+3/2 t^2

% v'=fy_1/W=-((6t^2+6t)e^(2t))/2=-(3t^2+3t)e^(2t) ----> v={-(3t^2+3t)/2  +(6t+3)/4 -6/8}e^(2t)=-3/2 t^2 e^(2t)

% y_p=(t^3+3/2 t^2)e^(t)-3/2 t^2 e^(t)=t^3e^t.  Yay!

% TABULAR INT BY PARTS:

%  u        | dv
%  -3t^2-3t | e^(2t)
%  -6t-3    | 0.5e^(2t)
%  -6       | 0.25e^(2t)
%   0       | 0.125e^(2t)





\item Using that $y_1=e^{t}$ and $\displaystyle y_2=\frac{e^t}{t}$ are a fundamental set of solutions to \\$ty^{\,\prime\prime}-2(t-1)y^{\,\prime}+(t-2)y=0$ on $(0,\infty)$ determine a possible choice for the function $v$ when using Variation of Parameters to find a particular solution of the form $y_p=uy_1+vy_2$ for $ty^{\,\prime\prime}-2(t-1)y^{\,\prime}+(t-2)y=e^{2t}$ on $(0,\infty)$.

    % f=e^(2t) /t.  W=e^t(e^t /t -e^t /t^2)-e^(2t) /t =e^(2t) /t -e^(2t) /t^2-e^(2t) /t= -e^(2t) /t^2

    % So u'=-fy_2/W=- (e^(2t) /t)(e^t /t) /(-e^(2t) /t^2) =e^t ---> u=e^t

    % v'=fy_1/W=(e^(2t) /t)(e^t)/(-e^(2t) /t^2)=-te^t ---> v=(1-t)e^t

    % y_p=(1/t)e^{2t}.  Sanity: y'=[(2/t)-(1/t^2)]e^{2t} and y''=[(4/t)-(4/t^2)+(2/t^3)] e^{2t}

    % t[(4/t)-(4/t^2)+(2/t^3)]-2(t-1)[(2/t)-(1/t^2)]+(t-2)(1/t) ]e^{2t}=

      % [4-(4/t)+(2/t^2)-4+(6/t)-(2/t^2)+1-2/t ]e^{2t}=

      % e^{2t}=YES


\begin{enumerate}[label=(\alph*)]
\item $v=t$
\item $v=-0.5e^{2t}$
\item $v=-e^{-t}$
\item $v=(t-1)e^t$
\item $v=-(t-1)e^t$ %***
\item $v=-e^t$
\end{enumerate}

%  (e) with options through (f)


\item Suppose $y(t)$ is the solution to $t^2 y\,^{\prime\prime}-2y=t^2$ on $(0,\infty)$ which \emph{involves the least number of terms}.   Determine $y(2)$.   Round to the nearest hundredth.

% 4*ln(2)/3\approx 0.92

% CP of CE:  (m+1)(m-2) ... y_1=t^2, y_2=1/t ...  W=-3 ... f=1.

% y_p=uy_1+vy_2 ... u'=1/(3t) and v'=-t^2/3 ... u=(1/3)ln(t) and v=-t^3/9

% y_p=(1/3)t^2ln(t)-t^2/9

% With least number of terms:  y=(1/3)t^2ln(t)


\item Suppose that $y=(1+t)^2$ is one solution to $y\,^{\prime\prime}+p(t)y\,^{\prime}+q(t)y=0$ on $(-\infty,\infty)$ where $p(t)$ and $q(t)$ are continuous functions on $(-\infty,\infty)$.  Given that the Wronskian $W(t)$ of any two solutions to $y\,^{\prime\prime}+p(t)y\,^{\prime}+q(t)y=0$ on $(-\infty,\infty)$ is constant, find the solution $y(t)$ to the IVP \\$y\,^{\prime\prime}+p(t)y\,^{\prime}+q(t)y=1+t$, $y(0)=0$, $y\,^{\prime}(0)=0$ on $(-\infty,\infty)$.  Determine $y(1)$.  Round to the nearest hundredth.

    % y(1)=17/24\approx 0.71



% Since W is constant, y_2=1/(1+t) can be used with y_1=(1+t)^2.  W=-3 and f=1+t

% u'=1/3 ... u=t/3

% v'=-(1+t)^3/3 ... v=-(1+t)^4/12

% y=t(1+t)^2/4+c_1(1+t)^2+c_2/(1+t)

% y(0)=0:  0=c_1+c_2

% y'=(1+t)^2/4+t(1+t)/2+2c_1(1+t)-c_2/(1+t)^2

% y'(0)=0:  0=1/4+2c_1-c_2 ...  0=1/4 +3c_1 ... c_1=-1/12 and c_2=1/12

% y=t(1+t)^2/4-(1+t)^2/12+1/[12(1+t)]


\item Suppose $y(x)$ satisfies $x^2y^{\,\prime\prime}-2y=-\dfrac{3}{x}$, $y(1)=0$, and $y^{\,\prime}(1)=1$ on $(0,\infty)$.  Determine $y(2026)$.  Round to the nearest thousandth.

% m(m-1)-2=m^2-m-2=(m+1)(m-2).

% y_1=1/x, y_2=x^2 and f=-3/x^3

% W=3.  So u'=-fy_2/W=1/x ---> u=ln(x)

% v'=fy_1/W=-1/x^4 ---> v=1/(3x^3)

% y_p=(1/x)ln(x)+1/(3x)

% Solution with least number of terms: y=ln(x)/x.   y'=1/x^2-ln(x)/x^2

% y(2026)=ln(2026)/2026 \approx 0.004

\item Given that $y_1=xe^x$ and $y_2=xe^{-x}$ form a fundamental set of solutions to $x^2y^{\,\prime\prime}-2xy^{\,\prime}-(x^2-2)y=0$ on $I=(0,\infty)$ find the particular solution $y_p$ to $x^2y^{\,\prime\prime}-2xy^{\,\prime}-(x^2-2)y=3x^4$ that involves the least number of terms.  Find $y_p(1)$.


% -3


    %  f=3x^2. W=x(1-x)-(1+x)x=-2x^2.  So u'=-3x^2 (xe^(-x))/(-2x^2)=1.5xe^(-x) ---> u=1.5(-xe^(-x)-e^(-x))

    % v'=3x^2 (xe^x)/(-2x^2)=-1.5xe^x ---> v=-1.5 (xe^x-e^x)

    % y_p=-1.5(x+1)x-1.5 (x-1)x=-1.5[(x+1)x+(x-1)x]=-3x^2  Sanity check.  y=-3x^2 so y'=-6x and y''=-6

    % -6x^2-2x(-6x)-(x^2-2)(-3x^2)=3x^4 YES!  So y_p(1)=-3


\item Suppose $y(x)$ is the particular solution to $2x^2 y^{\,\prime\prime}+xy^{\,\prime}-3y=x^{3/2}$ which \emph{involves the least number of terms}.  Determine $y(2026)$.  Round to the nearest tenth.

% (1/5) 2026^(3/2) ln(2026)\approx 138864.6

% y = (1/5)x^(3/2)ln(x)

\item Suppose $y(x)$ is the particular solution to $y^{\,\prime\prime}+4y=4\sec^2(2x)$ which \emph{involves the least number of terms}.  Determine $\displaystyle y\left(\frac{\pi}{6}\right)$.  Round to the nearest hundredth.

% y(pi/6) \approx 0.14

% y = -1+sin(2x) ln|sec(2x)+tan(2x)|

Hint: you may find $\displaystyle \int \sec (\theta)\, \mathrm{d}\theta=\ln \left|\sec (\theta)+\tan (\theta)\right|+C$ helpful.

\end{enumerate}



\end{document}
